\chapter{Formal Learning and Identification in the Limit}
\label{sec:p15-formal-identification}

The iid risk guarantees in \cref{sec:p15-vc-dimension-and-uniform-learning,
sec:p15-occam-learning} concern finite samples and errors weighted by
a sampling law. Identification in the limit instead asks whether one
learner eventually settles on an exactly correct hypothesis for every
allowed complete presentation. There is no probability of success
or announced sample-size cutoff in that criterion.
This chapter reconstructs the relevant original models and diagonals
of Gold's 1967 paper \cite{Gold1967IdentificationFLT}.

\section{The learner, its names and its information}
\label{sec:p15-flt-model}

\begin{hbdefinition}{def:p15-flt-names}{Languages and hypothesis indices}
Fix a finite nonempty alphabet \(A\).
Gold writes \(\Sigma A\) for its finite strings; write
\(\mathcal W=A^*\) for the same set, including the empty
word \(\varepsilon\). Fix a computable length-lexicographic
bijection \(b:\mathbb N\to\mathcal W\) with computable
inverse. It can be chosen primitive recursive on finite-string
codes, as required by the source's integer/string bridge.
A language is \(L\subseteq\mathcal W\).

Fix a naming map \(\Gamma\) from valid natural-number
hypothesis codes to languages. It may be partial and may
assign the same language to different codes.
A target class \(\mathcal C\) is contained in its range.
No uniformly terminating membership algorithm for
\(\Gamma(e)\) is assumed.
Gold's tester names are programs halting with membership
answer one or zero; his generator names are programs whose
produced range is the language. They are different input/output
interfaces, not interchangeable numerical indices.
\end{hbdefinition}

\begin{hbdefinition}{def:p15-flt-data}{Texts and informants}
A Gold text for a nonempty \(L\) is an infinite sequence
\(T(0),T(1),\ldots\) of words whose range is exactly \(L\).
Every positive word eventually occurs, in any order and with
arbitrary repetition; missing words are not labelled negative.
An informant is a sequence of pairs \((w,\ell)\), with
\(\ell=\mathbf1_L(w)\), that eventually presents every
word \(w\in\mathcal W\). Repetition must remain truthful.
A methodical informant presents \(b(n)\) at step \(n\).
In a request informant the learner chooses each word queried.

For empty languages and harmless waiting, the host also allows
a pause symbol \(\#\notin\mathcal W\) in texts, with the
nonpause range exactly \(L\).
This is an explicit extension of Gold's pure positive-word
text convention. The main diagonals below already use
nonempty languages and no pauses.
\end{hbdefinition}

\begin{hbdefinition}{def:p15-flt-ex}{Syntactic stabilization and semantic correctness}
An effective learner is a total computable function \(G\)
from finite data sequences to hypothesis codes, optionally
with an abstention symbol \(\bot\).
Gold's \(g_t=G(i_1,\ldots,i_t)\) is this same prefix rule.
It identifies \(L\) on presentation \(T\) if
\[
 (\exists e)(\exists n_0)(\forall n\geq n_0)\
 G(T{\upharpoonright}n)=e,\qquad \Gamma(e)=L.
\]
Learning a class quantifies over every \(L\in\mathcal C\)
and every admissible presentation of \(L\).
The index must stabilize and its meaning must be correct.
Stabilizing on a wrong language fails; alternating two
equivalent program names forever also fails, even if both
always describe the correct language.
The convergence stage may depend on the target and presentation.
\end{hbdefinition}

Gold separately considers arbitrary, recursive and primitive-recursive
texts, and tester versus generator naming
\cite[pp.\ 449--452]{Gold1967IdentificationFLT}.
Restrictions on the presenting sequence are not restrictions
on the target's membership algorithm.

\begin{hbtheorem}{thm:p15-flt-informant-channels}{proved=annex:p15-flt-informant-channels}{The three truthful informant channels}
Arbitrary fair, methodical and request informants give the same
identifiable classes, with computable finite-prefix learners.
A request strategy can be simulated from a fair informant by
caching answers and waiting for each requested word, repeating
the previous hypothesis while waiting.
The theorem concerns eventual identification, not equal time
or query complexity.
\end{hbtheorem}

\section{Positive results with actual effective enumerations}
\label{sec:p15-flt-positive}

\begin{hbdefinition}{def:p15-flt-uniform-family}{A uniformly decidable indexed family}
For this positive theorem, require a total computable
\(D:\mathbb N\times\mathcal W\to\{0,1\}\).
Let \(H_e=\{w:D(e,w)=1\}\).
This theorem uses family indices as names, \(\Gamma(e)=H_e\);
the fixed compiler transports them to Gold tester-program names.
Each member has an effectively obtained tester: hard-code
\(e\) into the program computing \(D(e,\cdot)\).
Duplicates are permitted. The joint totality hypothesis is
substantive; it is not inferred from the existence of
an individual decider for each \(H_e\).
\end{hbdefinition}

\begin{hbtheorem}{thm:p15-flt-enumeration}{proved=annex:p15-flt-enumeration}{Identification by enumeration from an informant}
Every such indexed family is identifiable from every truthful
fair informant. At data length \(n\), check indices
\(e\leq n\) against all observed labels and return the least
consistent one; if there is none, abstain.
Eventually the result is the least index for the target language.
\end{hbtheorem}

Gold's Theorem I.4 obtains this situation for primitive-recursive
languages. Enumerate their finite primitive-recursive descriptions;
the joint interpreter is total recursive, though it need not
itself be primitive recursive. Apply the source's Boolean map
\(w(z)=1\) exactly when \(z=1\), and zero otherwise.
It enumerates the primitive-recursive characteristic functions.
Finite structural evaluation and bounded primitive-recursion
loops justify termination on each description and input
\cite[pp.\ 467--468]{Gold1967IdentificationFLT}.
Thus the positive result includes that class, not all recursive
languages merely by listing their programs.

\begin{hbtheorem}{thm:p15-flt-finite-text}{proved=annex:p15-flt-finite-text}{All finite languages from positive text}
The class of finite languages is text-identifiable.
For a finite set \(F\), let
\[
 d(F)=\sum_{\{j:b(j)\in F\}}2^j.
\]
Its finite-language tester checks the corresponding bit of
\(d(F)\). A fixed effective compiler \(c(d)\) returns its
program index. The learner outputs \(c(d(\operatorname{content}
(\sigma)))\), where content discards pauses.
Its parameter eventually becomes the full target set and
then no longer changes. The pause extension includes the
empty language; Gold's pure-text result concerns its
nonempty presentations.
\end{hbtheorem}

\begin{hbexample}{def:p15-flt-trace}{An actual changing-hypothesis trace}
For the unary alphabet, take \(b(j)=a^j\) and target
\(\{\varepsilon,a^2\}\). A text can begin
\(\#,\ a^2,\ a^2,\ \varepsilon,\ \#,\ \varepsilon,\ldots\).
The finite-set parameters after those observations are
\(0,4,4,5,5,5,\ldots\), and the actual hypotheses are
\(c(0),c(4),c(4),c(5),c(5),c(5),\ldots\).
The numbers are set parameters, not invented raw program indices.
No finite pause proves that another word will never appear.
\end{hbexample}

\begin{hbtheorem}{thm:p15-flt-pr-text}{proved=annex:p15-flt-pr-text}{The restricted primitive-recursive-text generator result}
If presentations are required to be primitive-recursive
positive texts, the class of nonempty recursively enumerable
languages is generator-identifiable in the limit.
The learner identifies a primitive-recursive program for
the presenting sequence by enumeration and outputs its
generator index. It need not produce a membership tester.
Every nonempty r.e. language has such a presentation.
\end{hbtheorem}

This is Gold's Theorem I.7, the anomalous entry of his
Table I. It does not contradict arbitrary- or recursive-text
impossibility. It changes the permitted presentations and
the output naming relation together.

\section{The complete superfinite diagonal}
\label{sec:p15-flt-superfinite}

\begin{hbdefinition}{def:p15-flt-superfinite}{Superfinite classes}
A class is superfinite when it contains every finite
language and at least one infinite language.
This is a property of a class, not the statement that
one particular target has infinitely many words.
\end{hbdefinition}

\begin{hbtheorem}{thm:p15-flt-superfinite}{proved=annex:p15-flt-superfinite}{No positive-text learner for a superfinite class}
No learner identifies a superfinite class from every
arbitrary positive text, for any naming map covering the
class. The obstruction holds even without a computability
restriction on the learner.

For computable learners and Gold's generator naming, an
infinite r.e. target already has a recursive positive text
on which identification fails. Enumerate it without
repetition as \(a_1,a_2,\ldots\).
At stage \(r\), extend the existing prefix \(u\) by
\(a_{2r-1}^{\,x}a_{2r}^{\,y}\), searching positive
\(x,y\) with
\[
 G(u a_{2r-1}^{\,x})
 \ne G(u a_{2r-1}^{\,x}a_{2r}^{\,y}).
\]
The finite-target convergence assumptions ensure such a
pair exists at every stage. Its test is equality of
finite hypothesis codes, not semantic language equality.
\end{hbtheorem}

The annex supplies both the arbitrary-text construction and
this effective block diagonal, including fairness, termination
and computation of each text position.
Gold's Theorem I.8 is the recursive-text generator version
\cite[pp.\ 470--471]{Gold1967IdentificationFLT}.
Since the regular languages include every finite language
and an infinite one, positive-only unrestricted texts
cannot identify that whole class. The result does not say
that every smaller language class is unlearnable.

\section{Why all recursive hypotheses cannot be silently indexed}
\label{sec:p15-flt-index-boundary}

\begin{hbtheorem}{thm:p15-flt-no-uniform-list}{proved=annex:p15-flt-no-uniform-list}{No uniform total list of all recursive languages or functions}
There is no total computable joint decider \(D(e,w)\)
whose indexed languages include every recursive language.
Likewise no total computable \(p(e,n)\) lists every total
computable numerical function. Therefore the earlier
least-consistent algorithm cannot be extended to all
recursive objects by an unsupported enumeration assumption.
\end{hbtheorem}

\begin{hbtheorem}{thm:p15-flt-all-recursive}{proved=annex:p15-flt-all-recursive}{The stronger methodical-informant obstruction}
No computable learner identifies all recursive languages
even from methodical informants, with generator names.
Assuming it learns all finite and cofinite targets,
one can construct a recursive characteristic sequence by
successively finding \(x,y\geq1\) with
\[
 G(u0^x)\ne G(u0^x1^y)
\]
and appending those blocks. Each search terminates and
forces another change on the resulting recursive target.
\end{hbtheorem}

This reconstructs Gold's Theorem I.5
\cite[pp.\ 468--469]{Gold1967IdentificationFLT}.
It uses no universal decision algorithm for the guessed
languages. A tester learner could be compiled to a
generator learner by scanning words and emitting accepted
ones, so the negative result is not repaired by demanding
testers as output.

\begin{hbtheorem}{thm:p15-flt-no-certificate}{proved=annex:p15-flt-no-certificate}{Stabilization has no general finite certificate}
No procedure can always halt after finitely much positive
text and certify the exact target for all finite languages.
A prefix on which it halts for \(F\) can also begin a
text for \(F\cup\{z\}\) with a new word delayed until
after that prefix.
\end{hbtheorem}

Loss of fairness or truthfulness also invalidates the positive
proofs: an informant omitting the only distinguishing word
cannot eliminate a wrong smaller index, and erroneous labels
can eliminate the correct one permanently.
Fixed finite data may fit multiple hypotheses; the source's
promise is eventual exact identity along every permitted
infinite presentation, not identification at a known stage.

\begin{proofboundary}{def:p15-flt-scope}{Distinct learning criteria and source versions}
Syntactic stabilization, eventual semantic correctness,
finite certification, prediction accuracy and PAC risk
control are different criteria.
IID sampling does not generally give a fair text when some
target words have zero sampling probability. Conversely
fair presentations can delay a distinguishing word arbitrarily.
Neither framework's theorem supplies the other's quantifiers.
Gold's primitive-recursive-text tester obstruction, Theorem I.9,
also uses his separately cited limiting-recursion normalization;
it is not claimed to follow from the recursive block proof.
\end{proofboundary}
